\documentclass[11pt]{article}
\usepackage[T1]{fontenc}
\usepackage[utf8]{inputenc}
\usepackage{lmodern}
\usepackage[a4paper,margin=27mm,headheight=14pt]{geometry}
\usepackage{amsmath,amssymb,amsthm,mathtools}
\usepackage{microtype,booktabs,enumitem}
\usepackage[hidelinks]{hyperref}
\hypersetup{
 pdftitle={MI-27: The logarithmic commutator inequality with sharp constant one},
 pdfauthor={Sidney Holden},
 pdfsubject={Complete analytic proof, sharpness, and reproducible consistency checks},
 pdfkeywords={MI-27, logarithmic commutator, entropy, sharp constant, positive part}
}
\usepackage{fancyhdr}
\pagestyle{fancy}
\fancyhf{}
\fancyhead[L]{\small MI-27: sharp logarithmic commutator inequality}
\fancyhead[R]{\small 12 September 2026}
\fancyfoot[C]{\thepage}
\setlength{\parindent}{0pt}
\setlength{\parskip}{5pt plus 1pt}
\setlength{\emergencystretch}{2em}
\numberwithin{equation}{section}
\newtheorem{theorem}{Theorem}[section]
\newtheorem{lemma}[theorem]{Lemma}
\newtheorem{proposition}[theorem]{Proposition}
\newtheorem{corollary}[theorem]{Corollary}
\theoremstyle{remark}
\newtheorem{remark}[theorem]{Remark}
\DeclareMathOperator{\tr}{tr}
\DeclareMathOperator{\sgn}{sgn}
\DeclareMathOperator{\rank}{rank}
\newcommand{\Id}{I}
\newcommand{\one}[1]{\left\lVert #1\right\rVert_1}
\newcommand{\op}[1]{\left\lVert #1\right\rVert_\infty}
\newcommand{\ent}{\mathsf S}
\newcommand{\rel}{\mathsf D}
\newcommand{\hs}{\mathsf E}
\newcommand{\dd}{\,\mathrm d}
\newcommand{\ii}{\mathrm i}
\newcommand{\C}{\mathbb C}
\title{\textbf{MI-27: The logarithmic commutator inequality\\with sharp constant one}}
\author{Sidney Holden\\{\small Center for Computational Biology, Flatiron Institute}\\{\small Simons Foundation}}
\date{12 September 2026}

\begin{document}
\maketitle
\thispagestyle{empty}

\begin{abstract}
For every pair of positive definite complex matrices $A,B$ with
$\tr(A+B)=1$, we prove
\[
 \one{[B,\log(A+B)]}\le
 -\tr(A)\log\tr(A)-\tr(B)\log\tr(B).
\]
The proof uses Frenkel's integral representation of quantum relative
entropy, in its positive-part form given by Hirche and Tomamichel.
The key estimate is that $\tr(\rho-\gamma e^{\ii tH}\sigma e^{-\ii tH})_+$
is $\op{H}$-Lipschitz in $t$, uniformly for $\gamma\ge1$ and for all
matrix sizes. A positive integral of these estimates has total weight
exactly the binary entropy. A strictly positive definite $2\times2$
family from the supplied findings pack proves that no coefficient
smaller than one is possible.
\end{abstract}

\section{Statement and conventions}

The problem designated MI-27 in the linked collection asks specifically
for coefficient one, rather than merely a dimension-independent
coefficient \cite{MI27}. Its historical formulation and a coefficient-two
bound are recorded in \cite{AK,Audenaert}. The following settles the
stated matrix inequality.

\begin{theorem}[MI-27, with the optimal universal coefficient]
\label{thm:main}
Let $n\ge1$, and let $A,B\in\C^{n\times n}$ be positive definite with
$\tr(A+B)=1$. Put $a=\tr A$ and $b=\tr B$. Then
\begin{equation}
 \boxed{\one{[B,\log(A+B)]}\le h(b)},
 \qquad h(b)=-a\log a-b\log b,\quad a+b=1.
 \label{eq:main}
\end{equation}
The coefficient one is best possible even when $n=2$ and both matrices
are strictly positive definite.
\end{theorem}

All logarithms are natural. The commutator is $[X,Y]=XY-YX$;
$\one{X}=\tr\sqrt{X^*X}$ is the full trace norm, not one half of it;
and $\op{X}$ is the operator norm. A density matrix is a positive
semidefinite matrix of trace one. For Hermitian $X$, $X_+$ is its positive
part. Write
\[
 \ent(\rho)=-\tr\rho\log\rho,\qquad
 \rel(\rho\Vert\sigma)=\tr\rho(\log\rho-\log\sigma),
\]
with the usual support convention for relative entropy and $0\log0=0$.
Only positive definite matrices are needed in the application to
Theorem~\ref{thm:main}.

\textbf{Proof dependency.} The only non-elementary input is the published
relative-entropy identity in Section~\ref{sec:representation}.
Everything used after that identity is proved below. The argument does
not assume that $A$ and $B$ commute, that either has rank one, or that the
projection reduction in the supplied pack has already been solved.

\section{A uniform positive-part estimate}

\begin{lemma}[Commuting with an effect]
\label{lem:effect}
If $X\ge0$ and $0\le Q\le\Id$, then
\begin{equation}
 \one{[X,Q]}\le\tr X.
 \label{eq:effect}
\end{equation}
\end{lemma}
\begin{proof}
The Hermitian matrix $R=2Q-\Id$ is a contraction. Consequently, the
triangle inequality and the trace-norm product inequality give
\[
 \one{[X,Q]}=\tfrac12\one{XR-RX}
 \le\tfrac12\bigl(\one{XR}+\one{RX}\bigr)
 \le\one{X}\op{R}\le\tr X.
\]
\end{proof}

For density matrices $\rho,\sigma$ and $\gamma\ge1$, define
\begin{equation}
 \hs_\gamma(\rho\Vert\sigma)
   =\tr(\rho-\gamma\sigma)_+.
 \label{eq:hockey}
\end{equation}
This is often called the quantum hockey-stick divergence.

\begin{lemma}[Uniform unitary-orbit Lipschitz bound]
\label{lem:unitary}
Let $\rho,\sigma$ be density matrices, let $H=H^*$, and set
$U_t=e^{\ii tH}$ and $\sigma_t=U_t\sigma U_t^*$. For every $\gamma\ge1$
and every $s,t\in\mathbb R$,
\begin{align}
 |\hs_\gamma(\rho\Vert\sigma_t)
    -\hs_\gamma(\rho\Vert\sigma_s)|
     &\le |t-s|\op{H},\label{eq:orbit-second}\\
 |\hs_\gamma(\sigma_t\Vert\rho)
    -\hs_\gamma(\sigma_s\Vert\rho)|
     &\le |t-s|\op{H}.\label{eq:orbit-first}
\end{align}
\end{lemma}
\begin{proof}
For every Hermitian $M$, diagonalization shows that
\begin{equation}
 \tr M_+=\max_{0\le Q\le\Id}\tr QM.
 \label{eq:variational}
\end{equation}
In particular $|\tr M_+-\tr N_+|\le\one{M-N}$. Thus
$f(t)=\tr M(t)_+$, where $M(t)=\rho-\gamma\sigma_t$, is locally
Lipschitz, hence absolutely continuous on compact intervals and
differentiable almost everywhere.

At a point of differentiability, fix the spectral projection
$Q_t=\boldsymbol 1_{(0,\infty)}(M(t))$. It maximizes
\eqref{eq:variational}, and it commutes with $M(t)$. We have
\[
 f(t+u)\ge\tr Q_tM(t+u)
       =f(t)+u\tr Q_tM'(t)+o(|u|).
\]
The positive and negative choices of $u$, together with differentiability
of $f$ at $t$, imply
$f'(t)=\tr Q_tM'(t)$. No differentiability of $Q_t$ is being assumed.

Since $[\rho-\gamma\sigma_t,Q_t]=0$, we have
$[\rho,Q_t]=\gamma[\sigma_t,Q_t]$. Therefore
\begin{align}
 f'(t)
  &=-\ii\gamma\tr Q_t[H,\sigma_t]
    =-\ii\gamma\tr H[\sigma_t,Q_t]\notag\\
  &=-\ii\tr H[\rho,Q_t].\label{eq:key-cancellation}
\end{align}
Trace-norm duality and Lemma~\ref{lem:effect} now yield
\[
 |f'(t)|\le\op{H}\one{[\rho,Q_t]}\le\op{H}.
\]
Integrating this almost-everywhere bound for the absolutely continuous
function $f$ proves \eqref{eq:orbit-second}.

For \eqref{eq:orbit-first}, use unitary invariance:
\[
 \hs_\gamma(\sigma_t\Vert\rho)
   =\hs_\gamma(\sigma\Vert U_t^*\rho U_t),
\]
and apply the first result with generator $-H$.
\end{proof}

\begin{remark}
The cancellation in \eqref{eq:key-cancellation} removes the apparently
large factor $\gamma$. A generic perturbation bound on the second
argument would keep that factor and is not sufficient here. Eigenvalue
crossings at zero cause no gap: the scalar function is absolutely
continuous, and its derivative is used only almost everywhere.
\end{remark}

\section{The relative-entropy representation and its kernels}
\label{sec:representation}

Frenkel's Theorem~6 \cite{Frenkel}, in the equivalent form of
Hirche--Tomamichel, Corollary~2.3, equation~(2.22) \cite{HT}, states that
for density matrices $\rho,\sigma$,
\begin{equation}
 \rel(\rho\Vert\sigma)
  =\int_1^\infty
    \left\{\frac{\hs_\gamma(\rho\Vert\sigma)}{\gamma}
           +\frac{\hs_\gamma(\sigma\Vert\rho)}{\gamma^2}
    \right\}\dd\gamma.
 \label{eq:frenkel}
\end{equation}
This is the ordinary Umegaki relative entropy even when the matrices do
not commute. It is not a replacement of relative entropy by a different
quantum divergence.

For reference, Frenkel's original state-normalized form is
\[
 \rel(\rho\Vert\sigma)=
 \int_{-\infty}^{\infty}
 \frac{\tr\bigl((1-v)\rho+v\sigma\bigr)_-}{|v|(v-1)^2}\dd v.
\]
The integrand vanishes for $0\le v\le1$. In the part $v>1$, the
substitution $\gamma=v/(v-1)$ gives the first integral in
\eqref{eq:frenkel}. In the part $v<0$, the substitution
$\gamma=(v-1)/v$ gives the second. Here $X_-=(-X)_+$.

\begin{lemma}[Skew relative-entropy kernel]
\label{lem:skew}
For $0<\alpha<1$, put $\beta=1-\alpha$. Then
\begin{align}
 &\rel\bigl(\rho\Vert\alpha\rho+\beta\sigma\bigr)\notag\\
 &\quad=\int_1^\infty\left\{
 \frac{\beta^2}{\gamma(\beta+\alpha\gamma)^2}
       \hs_\gamma(\rho\Vert\sigma)
 +\frac{\beta^2}{(\alpha+\beta\gamma)^2}
       \hs_\gamma(\sigma\Vert\rho)
 \right\}\dd\gamma.
 \label{eq:skew-kernel}
\end{align}
\end{lemma}
\begin{proof}
Let $M=\alpha\rho+\beta\sigma$. For $1\le v<1/\alpha$,
\[
 \hs_v(\rho\Vert M)
   =(1-\alpha v)\,
     \hs_{\beta v/(1-\alpha v)}(\rho\Vert\sigma),
\]
and for $v\ge1/\alpha$ the positive part is zero. For all $v\ge1$,
\[
 \hs_v(M\Vert\rho)
   =\beta\,\hs_{(v-\alpha)/\beta}(\sigma\Vert\rho).
\]
Apply \eqref{eq:frenkel} to $(\rho,M)$. In its first integral make the
change of variables
\[
 \gamma=\frac{\beta v}{1-\alpha v},\qquad
 v=\frac{\gamma}{\beta+\alpha\gamma},\qquad
 \dd v=\frac{\beta}{(\beta+\alpha\gamma)^2}\dd\gamma.
\]
In its second integral put $v=\alpha+\beta\gamma$.
These give precisely the two coefficients in \eqref{eq:skew-kernel}.
Every integrand is nonnegative; the substitutions are valid also for
improper integrals.
\end{proof}

For $a,b>0$ with $a+b=1$, define the weighted entropy difference
\begin{align}
 \chi_{a,b}(\rho,\sigma)
   &:=\ent(a\rho+b\sigma)-a\ent(\rho)-b\ent(\sigma)\notag\\
   &=a\rel(\rho\Vert a\rho+b\sigma)
     +b\rel(\sigma\Vert a\rho+b\sigma).
 \label{eq:holevo}
\end{align}
The second equality follows by expanding the traces.

\begin{proposition}[A positive kernel of mass $h(b)$]
\label{prop:kernel}
For density matrices $\rho,\sigma$,
\begin{equation}
 \boxed{\displaystyle
 \chi_{a,b}(\rho,\sigma)
 =ab\int_1^\infty\left\{
 \frac{\hs_\gamma(\rho\Vert\sigma)}{\gamma(b+a\gamma)}
 +\frac{\hs_\gamma(\sigma\Vert\rho)}{\gamma(a+b\gamma)}
 \right\}\dd\gamma.}
 \label{eq:chi-kernel}
\end{equation}
The two nonnegative scalar kernels have respective masses
\begin{align}
 ab\int_1^\infty\frac{\dd\gamma}{\gamma(b+a\gamma)}
    &=-a\log a,\label{eq:mass-a}\\
 ab\int_1^\infty\frac{\dd\gamma}{\gamma(a+b\gamma)}
    &=-b\log b.\label{eq:mass-b}
\end{align}
\end{proposition}
\begin{proof}
Use Lemma~\ref{lem:skew} first with $(\alpha,\rho,\sigma)=(a,\rho,\sigma)$,
then with $(\alpha,\rho,\sigma)=(b,\sigma,\rho)$, and insert the results
in \eqref{eq:holevo}. The coefficient of
$\hs_\gamma(\rho\Vert\sigma)$ simplifies as follows:
\[
 \frac{ab^2}{\gamma(b+a\gamma)^2}
 +\frac{ba^2}{(b+a\gamma)^2}
 =\frac{ab}{\gamma(b+a\gamma)}.
\]
The other coefficient follows by exchanging $a$ and $b$.

For the first mass, an antiderivative is
$a\log\bigl(\gamma/(b+a\gamma)\bigr)$. At $\gamma=1$ it equals zero,
and at infinity it tends to $-a\log a$. The other mass is identical
with $a,b$ exchanged. Also $0\le\hs_\gamma\le1$ for density matrices
and $\gamma\ge1$, so \eqref{eq:chi-kernel} is absolutely convergent.
\end{proof}

\section{Proof of the upper bound}

\begin{proposition}[Finite-time entropy estimate]
\label{prop:finite}
Let $A,B>0$ with $\tr A=a$, $\tr B=b$, $a+b=1$, and let $H=H^*$.
For $S_t=A+e^{\ii tH}Be^{-\ii tH}$,
\begin{equation}
 |\ent(S_t)-\ent(S_0)|\le |t|\op{H}\,h(b)
 \quad\text{for all }t\in\mathbb R.
 \label{eq:finite}
\end{equation}
\end{proposition}
\begin{proof}
Set $\rho=A/a$, $\sigma=B/b$, and $\sigma_t=e^{\ii tH}\sigma e^{-\ii tH}$.
Unitary invariance of entropy gives
\[
 \ent(S_t)-\ent(S_0)
  =\chi_{a,b}(\rho,\sigma_t)-\chi_{a,b}(\rho,\sigma).
\]
Apply the two bounds of Lemma~\ref{lem:unitary} to the two integrands in
\eqref{eq:chi-kernel}. Integrate the resulting finite-difference bounds
and use \eqref{eq:mass-a}--\eqref{eq:mass-b}:
\begin{align*}
 |\ent(S_t)-\ent(S_0)|
 &\le |t|\op{H}\,ab\int_1^\infty
 \left\{\frac1{\gamma(b+a\gamma)}
          +\frac1{\gamma(a+b\gamma)}\right\}\dd\gamma\\
 &=|t|\op{H}(-a\log a-b\log b).
\end{align*}
The kernels are integrable and independent of $t$. In particular,
\emph{no differentiation under an improper integral is needed}.
\end{proof}

\begin{proof}[Proof of the inequality in Theorem~\ref{thm:main}]
Put $S=A+B$. The path $S_t$ in Proposition~\ref{prop:finite} is positive
definite near zero (in fact for every $t$), so its entropy is
differentiable. The derivative of the trace of a matrix function gives
\[
 \left.\frac{\mathrm d}{\mathrm dt}\ent(S_t)\right|_{t=0}
 =-\tr\bigl(S'_0(\log S+\Id)\bigr)
 =-\ii\tr[H,B]\log S
 =-\ii\tr H[B,\log S],
\]
where $S'_0=\ii[H,B]$ and $\tr S'_0=0$. Equivalently, the trace derivative
formula follows by diagonalizing $S$ and summing the diagonal terms of
the Fr\'echet derivative of $-x\log x$.

Divide \eqref{eq:finite} by $|t|$ and let $t\to0$. With the Hermitian
matrix $K=-\ii[B,\log S]$, this yields
\begin{equation}
 |\tr HK|\le\op{H}h(b)\qquad(H=H^*).
 \label{eq:dual}
\end{equation}
Choose $H=\sgn(K)$, using zero on $\ker K$. Then $\op{H}\le1$ and
$\tr HK=\tr|K|=\one{K}=\one{[B,\log S]}$. Inequality
\eqref{eq:dual} is exactly \eqref{eq:main}.
\end{proof}

\section{Sharpness within the strictly positive definite class}
\label{sec:sharp}

The following family is retained, with attribution, from the MI-27
findings in the user's supplied pack \cite{Pack}; its properties are
verified here rather than assumed.
For $0<t<1/2$, set
\begin{equation}
 S_t=\begin{pmatrix}1-t&0\\0&t\end{pmatrix},\qquad
 d_t=(1-2t^2)t(1-t),\qquad
 B_t=t^2S_t+d_t\begin{pmatrix}1&1\\1&1\end{pmatrix},\qquad
 A_t=S_t-B_t.
 \label{eq:family}
\end{equation}
Let $u_t=(\sqrt t,\sqrt{1-t})^{\mathsf T}$ and $P_t=u_tu_t^*$.
This is a rank-one orthogonal projection, and
\[
 S_t^{1/2}P_tS_t^{1/2}
   =t(1-t)\begin{pmatrix}1&1\\1&1\end{pmatrix}.
\]
Consequently,
\[
 B_t=S_t^{1/2}\bigl(t^2\Id+(1-2t^2)P_t\bigr)S_t^{1/2}.
\]
The middle matrix has eigenvalues $t^2$ and $1-t^2$, both strictly
between zero and one. Thus $A_t,B_t>0$, and
$\tr(A_t+B_t)=\tr S_t=1$. As an additional algebraic check,
\[
 \det A_t=\det B_t=t^3(1-t)(1-t^2)>0.
\]

Writing $b_t=\tr B_t$, direct calculation gives
\begin{equation}
 b_t=t^2+2d_t,\qquad
 \one{[B_t,\log S_t]}=2d_t\log\frac{1-t}{t}.
 \label{eq:family-values}
\end{equation}
Indeed, the commutator is real skew-symmetric with off-diagonal
magnitude $d_t\log((1-t)/t)$, so its two singular values both equal
that magnitude.

As $t\downarrow0$, $d_t/t\to1$ and $b_t/(2t)\to1$, whence
\[
 \frac{2d_t\log((1-t)/t)}{2t\log(1/t)}\longrightarrow1.
\]
Moreover, $h(x)=x\log(1/x)+x+O(x^2)$ as $x\downarrow0$.
Since $\log(1/b_t)/\log(1/t)\to1$, it follows that
\[
 \frac{h(b_t)}{2t\log(1/t)}\longrightarrow1.
\]
Therefore
\begin{equation}
 \boxed{\displaystyle
 \lim_{t\downarrow0}
 \frac{\one{[B_t,\log(A_t+B_t)]}}{h(\tr B_t)}=1.}
 \label{eq:sharp}
\end{equation}
For every $c<1$, the ratio in \eqref{eq:sharp} exceeds $c$ for sufficiently
small positive $t$. Thus no smaller universal coefficient is valid
under the original strict hypotheses. This completes the sharpness
part of Theorem~\ref{thm:main}.

\appendix
\section{The corresponding skew-divergence improvement}

Define the normalized skew divergence, for $0<\alpha<1$, by
\[
 \mathsf{SD}_\alpha(\rho\Vert\sigma)=
 \frac{\rel(\rho\Vert\alpha\rho+(1-\alpha)\sigma)}{-\log\alpha}.
\]
The kernels in \eqref{eq:skew-kernel} have masses
\begin{align*}
 \int_1^\infty\frac{\beta^2}{\gamma(\beta+\alpha\gamma)^2}\dd\gamma
   &=-\log\alpha-\beta,\\
 \int_1^\infty\frac{\beta^2}{(\alpha+\beta\gamma)^2}\dd\gamma
   &=\beta,
 \qquad\beta=1-\alpha.
\end{align*}
For the first integral one may use the antiderivative
$\log\gamma-\log(\beta+\alpha\gamma)+\beta/(\beta+\alpha\gamma)$;
for the second use $-\beta/(\alpha+\beta\gamma)$.
Thus Lemma~\ref{lem:unitary} also proves
\begin{equation}
 |\mathsf{SD}_\alpha(\rho\Vert U_t\sigma U_t^*)
    -\mathsf{SD}_\alpha(\rho\Vert\sigma)|
 \le |t|\op{H}.
 \label{eq:skew-Lipschitz}
\end{equation}
The same holds when the first argument is rotated, by joint unitary
invariance. This supplies the coefficient-one estimate in the
skew-divergence approach to mixing; compare the coefficient-two
unitary perturbation estimate in \cite[Section~8]{Audenaert}.
Indeed, expanding
\[
 \chi_{a,b}(\rho,\sigma)
 =-a\log a\,\mathsf{SD}_a(\rho\Vert\sigma)
  -b\log b\,\mathsf{SD}_b(\sigma\Vert\rho)
\]
and using \eqref{eq:skew-Lipschitz} gives a second presentation of
Proposition~\ref{prop:finite}.

\section{Connection to the pack's projection reduction}

The supplied pack \cite{Pack} reduced the target to the assertion
\begin{equation}
 \one{[S^{1/2}PS^{1/2},\log S]}
 \le h(\tr SP),\qquad S>0,\quad\tr S=1,
 \label{eq:projection}
\end{equation}
for every orthogonal projection $P$, with no rank restriction.
Theorem~\ref{thm:main} now proves this assertion: put
\[
 K_\varepsilon=\varepsilon\Id+(1-2\varepsilon)P,\quad
 B_\varepsilon=S^{1/2}K_\varepsilon S^{1/2},\quad
 A_\varepsilon=S-B_\varepsilon,
\]
apply the theorem for $0<\varepsilon<1/2$, and let $\varepsilon\downarrow0$.
The sum $S$ and its logarithm stay fixed, and all other quantities are
continuous. For $P=0$ or $P=\Id$, both sides of \eqref{eq:projection}
are zero, with $h(0)=h(1)=0$.

For completeness, the converse reduction is also valid. If $0\le K\le\Id$,
write its eigenvalues as $k_1\ge\cdots\ge k_n\ge0$, put $k_{n+1}=0$,
and let $P_j$ project onto the first $j$ eigenvectors. Then
\[
 K=\sum_{j=1}^n(k_j-k_{j+1})P_j+(1-k_1)0
\]
is a convex combination of projections. For the corresponding weights
$\theta_j$, put $B_j=S^{1/2}P_jS^{1/2}$ and $b_j=\tr B_j$.
Linearity, trace-norm convexity, \eqref{eq:projection}, and concavity of
$h$ give
\[
 \one{[S^{1/2}KS^{1/2},\log S]}
 \le\sum_j\theta_j h(b_j)
 \le h\!\left(\sum_j\theta_j b_j\right).
\]
Take $K=S^{-1/2}BS^{-1/2}$ to recover the original assertion.
Importantly, the $b_j$ are allowed to vary; no assertion is made about
extreme points in a fixed-trace slice. The upper-bound proof in the
main text does not require this reduction.

\section{Verification, scope, and provenance}

The accompanying \texttt{verification/verify.py} checks the rational
kernel identities and the sharpness-family determinants symbolically.
It also compares the relative-entropy and weighted-entropy integrals
with direct matrix evaluations, exercises a zero-eigenvalue crossing,
checks trace-norm duality and entropy derivatives, and evaluates the
strictly positive definite sharpness family at high precision. The
quadrature is split at generalized eigenvalues of the matrix pencils
to handle the positive-part kinks. The complete numerical output,
software versions, random seed, and tolerances are included.

These computations are consistency checks, not the proof of a universal
inequality. The analytic proof is Sections~1--5, using the published
identity \eqref{eq:frenkel}. The report is not a proof-assistant
formalization and has not been independently refereed. No claim of
historical publication priority is made. The source and exact input
pack provenance are recorded in \texttt{provenance.md}; the sharpness
family and projection reduction are explicitly credited to that pack.

\begin{thebibliography}{9}
\raggedright
\bibitem{MI27}
\emph{OpenProblemsInNLA}, ``MI-27: Constant one in the logarithmic
commutator inequality,'' matrix-inequalities-and-norms collection.
Accessed 12 September 2026.
\url{https://github.com/ajt60gaibb/OpenProblemsInNLA/tree/main/matrix-inequalities-and-norms/MI-27}.

\bibitem{AK}
K.~M.~R.~Audenaert and F.~Kittaneh,
\emph{Problems and Conjectures in Matrix and Operator Inequalities},
arXiv:1201.5232v3 (2012), Section~6, Conjecture~4, equation~(26),
and the following coefficient-one proposal.
\url{https://arxiv.org/abs/1201.5232}.

\bibitem{Audenaert}
K.~M.~R.~Audenaert,
\emph{Quantum Skew Divergence},
Journal of Mathematical Physics \textbf{55}, 112202 (2014), Section~8.
\url{https://arxiv.org/abs/1304.5935}.

\bibitem{Frenkel}
P.~E.~Frenkel,
\emph{Integral formula for quantum relative entropy implies data
processing inequality},
Quantum \textbf{7}, 1102 (2023), Theorem~6, pp.~6--7 of arXiv:2208.12194v4.
DOI: \href{https://doi.org/10.22331/q-2023-09-07-1102}{10.22331/q-2023-09-07-1102}.
\url{https://arxiv.org/abs/2208.12194}.

\bibitem{HT}
C.~Hirche and M.~Tomamichel,
\emph{Quantum R\'enyi and $f$-divergences from integral representations},
Communications in Mathematical Physics \textbf{405}, 208 (2024).
Corollary~2.3, equation~(2.22), p.~7 of arXiv:2306.12343v3;
see also Proposition~2.9 for skew-divergence kernels.
DOI: \href{https://doi.org/10.1007/s00220-024-05087-3}{10.1007/s00220-024-05087-3}.
\url{https://arxiv.org/abs/2306.12343}.

\bibitem{Pack}
User-supplied findings archive,
\emph{OpenProblemsInNLA matrix proof checkpoint, 2026-09-12},
MI-27/\texttt{result.md}, Theorems~1--2.
The archive supplies the projection equivalence and the strictly
positive definite sharpness family, but explicitly leaves the universal
upper bound unproved. An unchanged copy of that MI-27 result is included
in \texttt{provenance/input\_MI27\_result.md}.
\end{thebibliography}

\end{document}
